% ECONOMICS_PROBLEM: {"id": "C73-588", "title": "Strategic adaptation after regime changes", "jel_code": "C73", "source_id": "OP-0283"}
% !TeX program = xelatex
\documentclass[11pt,a4paper]{article}
\usepackage[margin=23mm]{geometry}
\usepackage{fontspec}
\setmainfont{Latin Modern Roman}
\usepackage{amsmath,amssymb,mathtools}
\usepackage{xcolor,enumitem,booktabs,longtable,array,xurl,hyperref}
\definecolor{muted}{HTML}{53636C}
\hypersetup{colorlinks=true,urlcolor=blue,linkcolor=blue}
\setlength{\parindent}{0pt}
\setlength{\parskip}{5pt}
\newcommand{\R}{\mathbb R}
\newcommand{\E}{\mathbb E}
\newcommand{\N}{\mathbb N}
\newcommand{\ind}{\mathbf 1}
\newcommand{\Var}{\operatorname{Var}}
\newcommand{\Cov}{\operatorname{Cov}}
\newcommand{\supp}{\operatorname{supp}}
\DeclareMathOperator*{\argmin}{arg\,min}
\DeclareMathOperator*{\argmax}{arg\,max}
\newcommand{\entrymeta}[3]{{\sffamily\small\color{muted}\textbf{\texttt{#1}}\quad|\quad #2\par #3\par}}
\newcommand{\Problemref}[1]{\texttt{#1}}
\newcommand{\JELref}[2]{\texttt{#2} (JEL \texttt{#1})}
\begin{document}
\section*{Strategic adaptation after regime changes}
\textbf{Problem ID:} \texttt{C73-588}\quad \textbf{JEL:} \texttt{C73}\par
% BEGIN ECONOMICS STATEMENT
\label{problem:OP-0283}
{\sffamily\small\color{muted}Original subject group: Game theory and strategic interaction\par}
\label{definitions:problem:30007565}
\textbf{Audit classification:} Published research agenda. Fudenberg–Levine p.153 and conclusion discusses reevaluation and regime changes. The latent two-model experiment is editorial.
\subsection*{Definitions and precise research target}
\textbf{Model and research target.} Let human players repeatedly interact in a finite game with known action sets. At an experimentally randomized date \(\tau\), an institution changes the payoff matrix from \(u^0\) to \(u^1\). Players may receive an announcement or learn the change only through realized outcomes; this disclosure \(Z\) is randomized by interacting group. A candidate learning model maintains latent state \(M_{it}\in\{\text{old model},\text{revised model}\}\) and parameter estimates conditional on that state. Observe actions, realized payoffs, and incentivized predictive beliefs \(b_{it}\). Define revision time \(R_i\) by a prespecified belief-consistency criterion, rather than by an unobservable analyst judgment. The principal estimands are
\[
 \Pr(R_i-\tau\leq h\mid do(Z))
 \quad\text{and}\quad
 E\left[\sum_{t=\tau}^{\tau+h}(u_i^{\mathrm{best}}(t)-u_i(a_t))\mid do(Z)\right],
\]
where \(h\) is a fixed horizon and \(u_i^{\mathrm{best}}(t)=\max_{a_i}u_i^1(a_i,a_{-i,t})\). Determine whether observations distinguish changing beliefs within one model from switching the model itself, and derive identified sets when they cannot. Estimate disclosure effects and validate revision triggers on new changes. Regret here is a finite experimental outcome, not an equilibrium-existence target. Group randomization and explicit signal measurement assumptions are necessary because opponents also adapt after the change.

\emph{Definitions and maintained assumptions (editorial unless a source definition is named).} Fix finite actions, both payoff tensors, randomized change date \(\tau\) and disclosure treatment \(Z\), and whether \(\tau\) is known. At each date record the whole action profile and each player’s prediction distribution over declared outcomes. The old-model and revised-model hypotheses must specify observable prediction laws; labels alone do not define statistically distinguishable latent states. For tolerance \(\eta\) and window \(w\), define \(R_i=\inf\{t\ge\tau:\|b_{is}-b^{\rm revised}_{is}\|\le\eta\text{ for }s=t,\ldots,t+w-1\}\), with \(+\infty\) if never satisfied; the benchmark forecast and norm are part of the protocol. The sum of realized best-response shortfalls compares the chosen action with the best action against opponents’ realized actions, not a feasible counterfactual policy that changes their behavior. Censoring at experimental termination is reported explicitly. Randomization is by interacting group, and measurements before disclosure cannot be conditioned on treatment-induced post-change variables without changing the estimand.
\subsection*{Variants and assumption changes}
\begin{enumerate}
\item \textbf{Known change (editorial).} Announce both payoff tensors and the date. Estimate response speed to information and distinguish it from slow payoff learning.
\item \textbf{Hidden change (editorial).} Hide the changed tensor and provide noisy payoff feedback. Characterize observational equivalence between continuous parameter updating and discrete model replacement using additional randomized diagnostic outcomes.
\end{enumerate}
\subsection*{References and source audit}
\begin{enumerate}
\item p.153 and conclusion pp.165–166. Supports the stated research area; exact displayed specification is editorial. \url{https://economics.mit.edu/sites/default/files/2022-09/Whither%20Game%20Theory%20Towards%20a%20Theory%20of%20Learning%20in%20Games.pdf}
\end{enumerate}
\begin{enumerate}\raggedright
\item\hypertarget{definitions:source:30007565:1}{} Drew Fudenberg; David K. Levine. 2016. \emph{Whither Game Theory? Towards a Theory of Learning in Games}. Introduction, p.153; conclusion, pp.165–166.
\url{https://economics.mit.edu/sites/default/files/2022-09/Whither%20Game%20Theory%20Towards%20a%20Theory%20of%20Learning%20in%20Games.pdf}.
DOI: \url{https://doi.org/10.1257/jep.30.4.151}.
\end{enumerate}

\subsection*{Common conventions: Source mathematical conventions}

These conventions apply to each entry unless explicitly replaced.

\textbf{Models and identification.} A primitive model \(m\) specifies all probability laws, feasible choices, information, equilibrium and measurement rules used by its target. The admissible class \(\mathcal M\) is fixed independently of outcomes. If \(P_O(m)\) is its observable probability law and \(\tau(m)\) the target, its identified set at observable law \(P\) is
\[
 \mathcal I_\tau(P)=\{\tau(m):m\in\mathcal M,\ P_O(m)=P\}.
\]
Point identification means that this set is a singleton. An empty set signals incompatibility of data and assumptions. Coordinate bounds need not describe the joint set. Bounds are sharp only if every retained target is attainable or arbitrarily approachable by compatible models. Finite bounds require explicit restrictions on unobserved quantities; unrestricted confounding, hidden wealth, extrapolation or arbitrary equilibrium selection can yield uninformative sets.

\textbf{Causal interventions.} A potential outcome \(Y_i(p)\) is the outcome under a complete policy regime \(p\), including timing, financing, information and market spillovers. The target population and units are fixed before treatment. With interference, use the complete assignment vector or a justified exposure mapping; individual no-interference notation alone cannot identify an economy-wide intervention. A difference of mean potential outcomes does not require the cross-world joint distribution. Conditioning on post-treatment migration, survival, enrollment or employment changes the estimand and needs separate justification.

\textbf{Statistical statements.} An estimator, confidence set, forecast or test specifies a sampling law, model class and loss/coverage criterion. Uniform \((1-\alpha)\)-coverage means \(\inf_{m\in\mathcal M}\Pr_m\{\tau(m)\in C_n\}\geq1-\alpha\), or an explicitly asymptotic version. The whole parameter space trivially has coverage, so informativeness requires a stated width, risk or power objective. A forecasting improvement is relative to a named benchmark, information set and evaluation population.

\textbf{Equilibrium and optimization.} An equilibrium comprises feasible plans, optimality, beliefs where needed, budget constraints and all market/resource clearing. The primitives or a stated benchmark define each correspondence \(\mathcal E(p;m)\). Nonemptiness is assumed or proved, never inferred just from compact policy sets. A selected-equilibrium target uses a declared selection rule; a robust target quantifies over all permitted equilibria. Use suprema or infima unless attainment is established. An \(\varepsilon\)-optimal policy is feasible and within \(\varepsilon\) of the finite optimum value. Report an empty feasible set separately.

\textbf{Welfare and accounting.} Normative utility, interpersonal normalization, social weights, numeraire, discounting and the use of public receipts are primitives. Behavioral utility need not coincide with the welfare criterion. Household welfare includes ownership income and rents once; adding profits again to compensating variation generally double-counts them. A monetary equivalent is defined by an explicit transfer and price convention, with monotonicity and range conditions. Average effects, mechanism contributions, transfers and welfare losses are different targets. Infinite sums require convergence and dynamic feasible plans require terminal or no-Ponzi conditions.

\textbf{Source versions.} A working-paper date, revision date and journal publication year are distinguished. A DOI for a journal version is not automatically the DOI of an earlier manuscript. Locators refer to the stated version and printed pages where specified. A metadata-only check does not verify a claim about a full-text conclusion. Every entry's source subsection narrows the provenance claim accordingly.
% END ECONOMICS STATEMENT
\end{document}
